\subsection{ARCH-136-DISK-SUPPORT: uniform bounded torsion in the free horizontal source}
\label{sec:arch136disksupport}

This is a stronger partial-source conclusion, not integral
connectedness.  It concerns the uniformity required by cote 136,
physical pp.~117 and 119, footers 116 and 118.  The category and
coefficient algebra are those of AQ-136-COMPARISON, and the horizontal
cell attachment is the one explicitly constructed in
ALPOGE\_BRIDGES checkpoint 09, AQ-136-PARTIAL-SOURCE, proof, horizontal
step.  The extra relative $H^1$ repairs of that proof will not be used.

\paragraph{Statement.}
Let $C_1,\ldots,C_t\in\mathcal A_{\mathbf Z}$ be a finite nonempty
family.  For each form degree $q$, fix $\mu(q)\geq1$ so that every
row $(C_j)_*^{w,q}$ is horizontally acyclic for $w\geq\mu(q)$.
There is a simplicial strict bigraded differential PD algebra $H$
over $S=\Gamma_{\mathbf Z}(T)$, with strict maps $H\to C_j$, such
that
\[
 H_0=S,\qquad H_n^{0,*}=\mathbf Z[0],\qquad
 H_r(NH_*^{w,q})=0\quad(w\geq\mu(q),\ r\geq0),             \tag{136.DS1}
\]
and, in every weight and simplex,
\[
 H^0(H_n^{w,*})=S^w,
 \qquad
 L_d H^d(H_n^{w,*})=0\quad(d\geq1),\qquad
 L_d=\operatorname{lcm}\{1,\ldots,\lfloor d/2\rfloor\}.     \tag{136.DS2}
\]
The least common multiple of an empty set is $1$.
In particular the positive vertical cohomology vanishes through
degree three in every weight; the rational vertical cohomology is
exactly $S_{\mathbf Q}$ in every weight; and, for every prime $p$,
\[
 H^d((H_n)_{(p)}^{w,*})=0
                   \qquad(1\leq d\leq2p-1).               \tag{136.DS3}
\]
All these assertions are uniform in $n$ and $w$.  In a fixed degree
$d$, every defect is supported at the finitely many primes
$p\leq\lfloor d/2\rfloor$, with the displayed common annihilator.

\paragraph{Proof.}
Start with the constant simplicial algebra $S$.  Given a stage
$A$ with maps to all the $C_j$, attach a free horizontal cell for
every normalized horizontal cycle $z\in N_r A^{w,q}$ with
$w\geq\mu(q)$.  More explicitly, let $E_q$ be the vertical module
with generators $a,b$ in form degrees $q,q+1$, with $da=b$, and
let $\mathbb S^r\hookrightarrow\mathbb D^{r+1}$ be the Dold--Kan
image of the boundary inclusion from $\mathbf Z[r]$ into the
horizontal unit disk.  The cycle pair $(z,dz)$ defines the attaching
map from the free differential PD algebra on
$\mathbb S^r\otimes E_q$.  Push out along the free algebra on
$\mathbb D^{r+1}\otimes E_q$.
For each target choose a horizontal primitive of the image of $z$
in its form-$q$ row; its vertical derivative is the image of the
partner generator.  This extends every strict target map.  The
same free universal property checks all simplicial identities.
These are exactly the horizontal attachments of the cited supplied
proof; no vertical exactness of $A$ is required to perform them.

The boundary inclusion is an isomorphism in simplicial degree zero
when $r=0$, and both sides are zero there when $r>0$.  Thus every
stage has degree zero exactly $S$.  All added weights are positive.
At each fixed simplex the boundary inclusion has a free module
complement.  Tensor its splitting with $E_q$; the splitting then
commutes with the vertical differential.  The pushout is therefore,
as a vertical differential PD algebra, a free extension by
independent unit disks $da=b$.  Iterating this observation shows
that in every simplex the resulting algebra is isomorphic to
$S$ tensored with free vertical PD disks.  The choices of these
splittings need not be simplicial: they are used only to compute
vertical cohomology at each simplex, with bounds independent of
those choices.  There are no extensions with a differential
$da=p b$ among the free generator pairs.

Iterate the simultaneous attachment countably many times and take
the filtered colimit $H$.  Every old horizontal homology class in
the indicated weight tails is killed at the next stage.  Exactness
of filtered colimits of modules and exactness of normalization give
(136.DS1), with the original fixed cutoffs.  At a fixed simplex,
compatible splittings can be chosen at successive stages, so that
the colimit is still the tensor algebra of a possibly infinite
family of independent vertical unit disks.  Every monomial involves
only finitely many of those disks.

We compute this tensor product integrally.  If $q$ is even, its
one-disk algebra is $\Gamma(a)\otimes\Lambda(b)$.  Apart from its
constant, it is a direct sum of the two-term complexes
\[
 \mathbf Z a^{[r]}\xrightarrow{\ 1\ }
       \mathbf Z a^{[r-1]}b\qquad(r\geq1),
\]
in form degrees $qr,qr+1$.  These are integrally contractible,
including when $q=0$, by weight-preserving homotopies.
If $q$ is odd, its algebra is $\Lambda(a)\otimes\Gamma(b)$ and
its nonconstant blocks are
\[
 \mathbf Z a b^{[r-1]}\xrightarrow{\ r\ }
       \mathbf Z b^{[r]},
 \quad\text{in degrees }(q+1)r-1,(q+1)r,
 \quad r\geq1.                                           \tag{136.DS4}
\]
This follows from the special PD derivation identity and
$b b^{[r-1]}=r b^{[r]}$.  The block $r=1$ is contractible.
After removing these contractible summands and the even disks,
the full complex is a direct sum of tensor products of blocks
(136.DS4) with $q$ odd and $r\geq2$, together with its constant
summand $S$.  The decomposition preserves weight.  It applies to
infinitely many disks because each tensor monomial has finite
support.  Tensoring a contraction with any complex is legitimate
here; its usual signed homotopy is an actual contraction, not a
K\"unneth assumption.

A nonconstant component with just one block has cohomology only
in its upper degree $d=(q+1)r$, namely $\mathbf Z/r$.  Then
$r\leq\lfloor d/2\rfloor$, and $r$ divides $L_d$.
A component with $k\geq2$ blocks has its first term in degree
\[
             \ell=\sum_{i=1}^k((q_i+1)r_i-1).
\]
If it contributes in degree $d$, then $d\geq\ell$.  Every other
block starts in degree at least three, so for each $i$
\[
 2r_i-1+3(k-1)\leq d,
 \qquad r_i\leq\left\lfloor\frac{d-2}{2}\right\rfloor
                         \leq\lfloor d/2\rfloor.
\]
Multiplication by $r_i$ on its two-term block is chain homotopic
to zero: the homotopy sends the upper basis vector to the lower
basis vector.  Tensor this homotopy with the other factors, with
the usual signs.  Hence the cohomology of the component is killed
by $r_i$, and therefore by $L_d$.  Components with first term above
$d$ contribute nothing.  The zero-block component has no positive
cohomology.  Direct sums and the free coefficient factor $S$
preserve these conclusions, proving (136.DS2).

The same decomposition has $H^0=S$; all other blocks start in
positive degree and have injective lower differential over
$\mathbf Z$.  Rationally every multiplier $r$ is a unit, so the
whole augmentation ideal is contractible.  Finally localization
commutes with cohomology, and $p$ does not divide $L_d$ for
$d\leq2p-1$.  Thus (136.DS3) follows from (136.DS2).
This completes the proof.\hfill$\square$

\paragraph{Scope of the improvement.}
The supplied partial-source theorem left all higher vertical
cohomology as unspecified integer torsion.  For the pure horizontal
source the defects now have a uniform degree-dependent annihilator,
and the low vertical range improves from degrees zero and one to
zero through three without a relative repair.  This does not yet
make $H$ admissible: a nonzero fixed-degree torsion class may have
coefficient multiples in arbitrarily high weights.  The next
statement shows why even this stronger source cannot be repaired
by a source refinement preserving its low generators.

\subsection{ARCH-136-LOW-REFINEMENT: an obstruction for the inverse direction}
\label{sec:arch136lowrefinement}

\paragraph{Statement.}
For every integer $m\geq2$ there is a simplicial strict differential
PD algebra $A$ over $S$, with $A_0=S$, a strict map to the standard
admissible theory $B$, and
\[
 H_r(NA_*^{w,q})=0\quad(w>mq),
 \qquad H^0(A_n)=S,
 \qquad L_dH^d(A_n^{w,*})=0,                               \tag{136.LR1}
\]
such that no admissible $R$ admits an $S$-preserving strict map
$R\to A$ lifting a specified element of $A_1^{m,1}$.
Consequently there is no such source refinement that is surjective
in all form degrees below $2p-2$, for any prime $p$.
The example is simplexwise free on vertical unit disks; it has none
of the generator relations $dv=p u$ from AQ-136-HIGHER-COMPLETION.

\paragraph{Proof.}
Use the standard simplicial coordinate module $L$, with its
constant submodule generated by the coordinate sum, and put
$Q=L/\mathbf Z$.  Thus $NQ=\mathbf Z[1]$, $Q_0=0$, and
$Q_n\simeq\mathbf Z^n$.  Let $V,U$ be two copies of $Q$ in form
degrees one and two, both given weight $m$, and set
\[
 D=\Lambda(V)\otimes\Gamma(U),\quad dv=u,
 \qquad A=B\otimes_{\mathbf Z}D.                          \tag{136.LR2}
\]
The differential uses the identity between the two copies of $Q$;
the structure is therefore simplicial.  The coefficient map goes
into $B$.  Projection onto $B$ kills $V,U$ and is strict.  Since
$Q_0=0$, $A_0=S$.  In every simplex $A$ is $S$ tensored with the
even zero-form disks of $B$ and the odd unit disks of $D$, so
(136.DS2) proves its claimed vertical properties.

For the horizontal bound, a monomial in $D$ of form degree $d$
and weight $e$ satisfies $e\leq md$.  A form-$q$, weight-$w$
summand of $A$ is
\[
                  B^{w-e,q-d}\otimes D^{e,d}.
\]
If $w>mq$, then
$w-e>m(q-d)\geq q-d$.  The standard row of $B$ in this range is
simplicially contractible, by AQ-136-DP-MODEL.  Tensoring its linear
simplicial homotopy with the identity on the other factor proves
acyclicity, without a flatness inference.  For fixed $q$ the
possible pairs $(e,d)$ are finite; their direct sum gives the
first assertion of (136.LR1).

Choose the generator $a$ of $V_1\simeq\mathbf Z$, and write
$b=da$ for the generator of $U_1$.  The integral PD contraction
of $B_1$ onto $S$ identifies vertical cohomology of $A_1$ with
that of
\[
                     S\otimes\Lambda(a)\otimes\Gamma(b).
\]
For every prime $p$, the class of $b^{[p]}$ has exact order $p$
in degree $2p$ and weight $mp$, by the block (136.DS4).
Moreover every
\[
                    T^{[t]}b^{[p]}\qquad(t\geq0)           \tag{136.LR3}
\]
is nonzero of order $p$.  The coefficient variable $T$ is an
independent free tensor factor, so no binomial multiplier is
introduced in (136.LR3).

Suppose $f:R\to A$ were admissible in its source and some
$x\in R_1^{m,1}$ mapped to $a$.  For every $w\geq mp$, the element
\[
                 T^{[w-mp]}\gamma_p(dx)\in R_1^{w,2p}
\]
is closed.  Vertical admissibility makes it a boundary once
$w$ is sufficiently large.  Its image is the nonboundary
(136.LR3), a contradiction.  Thus $a$ cannot lift.  Since
$1<2p-2$ for $p\geq2$, preserving or surjecting onto all form
degrees below $2p-2$ would lift $a$, proving the consequence.
\hfill$\square$

\paragraph{Relation to the existing obstruction and the remaining direction.}
AQ-136-HIGHER-COMPLETION excludes certain extensions under a partial
source while retaining two specified target maps; its generators
have $dv=p u$.  The present statement concerns the opposite arrow
direction and has $dv=u$: an admissible source cannot refine this
free partial algebra while keeping its low generators.  It refutes
an inverse tower of such refinements if the limiting admissible
object still maps to the initial partial source and those low
generators lift.  It does not refute a forward sequence under the
partial source, a replacement changing the low-degree maps, or
integral common sources.  Indeed the example has a strict arrow
$B\to A$; that arrow omits the obstructing generator.
